% === Section 6: Meta-Analysis ===
\section{Meta-Analysis}
\label{sec:metaanalysis}\label{sec:meta-regression}\label{sec:metareg}

This section pools the extraction corpus (\S4) with fev-bench
\citep{shchur2025fevbench}, GIFT-Eval \citep{aksu2024gift}, and our
own gap-filling experiments (\S5) into a model-level meta-regression
of foundation models versus conventional baselines on retail demand
forecasting.  The restructured pipeline yields $k = 420$ paired
comparisons across 9~datasets and 11~foundation models, a substantial
increase from the $k = 10$ bucket-level analysis of the original
submission.

\subsection{Pooling strategy and effect size}
\label{subsec:pooling-strategy}

The unit of analysis is one paired comparison per (FM model, baseline,
task, metric) tuple.  For each within-paper group (identified by
\texttt{paper\_id}), we compute:

\begin{equation}
\Delta_{ij} \;=\; \text{metric}_{\text{FM}_i} - \text{metric}_{\text{baseline}_j}
\label{eq:delta-model-level}
\end{equation}

\noindent where $\text{baseline}_j$ is the within-group average of the
strongest non-FM baseline family: gradient-boosted trees (LightGBM,
CatBoost) where available (fev-bench, Source~B), or statistical models
(Auto~ETS, Seasonal Naive) where tree baselines are absent (GIFT-Eval).
The sign is oriented so that negative values mean ``FM beats
baseline''.

The variance proxy for each pair is $v_i = 1/n_{\text{FM}} +
1/n_{\text{baseline}}$, where $n$ is the number of series in the
evaluation set for each side.  For fev-bench tasks with known
$n_{\text{series}}$ (e.g., M5: 30{,}490; Favorita: 1{,}579), the
proxy is tight.  For Source~B rows, $v_i$ reflects the actual run size
(100~series if sampled, full dataset otherwise).

Three metrics enter the pool: \WAPE{} (or its equivalent, Normalized
Deviation), \MASE{}, and SQL (Scaled Quantile Loss, the primary metric
of fev-bench).  Rather than converting all to a single metric, we treat
\texttt{metric\_name} as a moderator (\S\ref{subsec:moderator-outcomes})
and report metric-specific sensitivity analyses.

\paragraph{Meta-regression model.}
The primary model is a multivariate random-effects specification with
dataset-level nesting:

\begin{alltt}
rma.mv(yi = delta, V = vi,
       random = ~ 1 | dataset_norm / paper_id,
       test = "t", method = "REML")
\end{alltt}

\noindent REML estimation with Knapp--Hartung small-sample degrees of
freedom \citep{knapp2003improved,viechtbauer2010metafor}.  The nesting
structure accounts for within-dataset clustering (multiple tasks from
the same dataset share confounds) and within-paper clustering (multiple
FM models from the same benchmark share evaluation protocol).

\paragraph{Data sources.}
The 420~pairs come from three sources:
\begin{itemize}
  \item \textbf{fev-bench} \citep{shchur2025fevbench}: 20 retail tasks
    $\times$ 7~FMs paired against LightGBM~(Recursive) and
    CatBoost~(Recursive) baselines, yielding 318~pairs across SQL,
    \MASE{}, and \WAPE{}.
  \item \textbf{GIFT-Eval} \citep{aksu2024gift}: 4~Sales-domain
    configurations (Car Parts, Hierarchical Sales $\times$\,2,
    Restaurant) $\times$ 7~FMs paired against Auto~ETS and Seasonal
    Naive baselines, yielding 84~pairs.
  \item \textbf{Source~B} (\S5): 9~cells (3~datasets $\times$
    3~horizons) $\times$ 4~FMs paired against the averaged baseline
    (\LGBM{} + Seasonal Naive) from the Azure sweep, yielding 36~pairs
    on \WAPE{}.
\end{itemize}

\subsection{Primary results}
\label{subsec:cross-paper-pooled}\label{sec:cross-paper-pooling}\label{sec:rq1-accuracy}

\paragraph{Full-pool intercept ($k = 420$).}

\begin{quote}
$\hat{\Delta}\;\text{(FM} - \text{baseline)} = -0.132$,\quad
$\text{SE} = 0.052$,\quad
$t(419) = -2.54$,\quad $p = 0.011$,\quad
95\,\% CI $[-0.234,\; -0.030]$.
\end{quote}

\noindent Foundation models are, on average, significantly better than
conventional baselines across all metrics, datasets, and model sizes in
this pool.  The between-dataset variance component is
$\hat{\sigma}^2_{\text{dataset}} = 0.012$; the within-dataset
between-paper component is
$\hat{\sigma}^2_{\text{paper}} = 0.031$.

\paragraph{Heterogeneity.}
$Q(419) = 115{,}285$, $p < .0001$.  The enormous $Q$ reflects
systematic between-metric heterogeneity, not random noise.  As shown in
the metric moderator analysis below, the direction and magnitude of the
FM advantage depends critically on which metric is used.

\subsection{The metric moderator: SQL versus \WAPE{}}
\label{subsec:metric-moderator}

The single most important finding of the meta-regression:

\begin{table}[htbp]
\centering
\caption{Metric-specific pooled $\hat{\Delta}$ from sensitivity analyses.
  Negative = FM better.}
\label{tab:metric-split}
\begin{tabular}{lrrlrl}
\toprule
Metric & $k$ & $\hat{\Delta}$ & 95\,\% CI & $p$ & Interpretation \\
\midrule
SQL (quantile loss)  & 128 & $-0.345$ & $[-0.442,\; -0.248]$ & $< .0001$ & FM strongly better \\
\MASE{}              & 128 & $-0.218$ & --- & --- & (intercept from moderator) \\
\WAPE{} / ND         & 164 & $+0.097$ & $[-0.086,\; +0.279]$ & $0.30$    & Null / FM marginally worse \\
\bottomrule
\end{tabular}
\end{table}

\noindent On SQL --- the primary metric of fev-bench, which measures the
full predictive distribution via quantile losses --- foundation models
are \textbf{strongly and significantly better} than tree baselines
($\hat{\Delta} = -0.345$, $p < .0001$).  On \WAPE{} --- a point
forecast accuracy metric --- the advantage \textbf{vanishes}
($\hat{\Delta} = +0.097$, $p = 0.30$, CI includes zero).

This metric split resolves the apparent contradiction in the literature:
papers that evaluate FMs on distributional metrics (CRPS, WQL, SQL)
report FM superiority; papers that evaluate on point-forecast metrics
(\WAPE{}, MAE) report parity or slight FM inferiority.  Both are
correct descriptions of the same underlying models.  The mechanism is
that FMs produce well-calibrated predictive distributions (trained on
probabilistic losses at scale) but their point forecasts (median or
mean of that distribution) are not systematically more accurate than
a well-tuned \LGBM{} point forecast.

\subsection{Forest plots}
\label{subsec:forest-plots}\label{sec:per-dataset-results}

Figures~\ref{fig:forest_m5}--\ref{fig:forest_car_parts} present
per-dataset forest plots of the model-level $\Delta$ pairs.  Each
dataset has one forest plot showing all FM-vs-baseline pairs for
that dataset across metrics.

\begin{figure}[htbp]
\centering
\includegraphics[width=\textwidth]{figure_6_forest_m5.pdf}
\caption{Forest plot for M5: model-level $\Delta$ across SQL, \MASE{},
  and \WAPE{}.  The FM advantage is large on SQL (left of zero) and
  reverses on \WAPE{} for several models.}
\label{fig:forest_m5}
\end{figure}

\begin{figure}[htbp]
\centering
\includegraphics[width=\textwidth]{figure_6_forest_favorita.pdf}
\caption{Forest plot for Favorita.  All FMs are close to or
  slightly better than baselines on all metrics.}
\label{fig:forest_favorita}
\end{figure}

\begin{figure}[htbp]
\centering
\includegraphics[width=\textwidth]{figure_6_forest_rossmann.pdf}
\caption{Forest plot for Rossmann.  Strong FM advantage across all
  metrics, driven by consistent SQL gains.}
\label{fig:forest_rossmann}
\end{figure}

\begin{figure}[htbp]
\centering
\includegraphics[width=\textwidth]{figure_6_forest_rohlik.pdf}
\caption{Forest plot for Rohlik~v2.}
\label{fig:forest_rohlik}
\end{figure}

\begin{figure}[htbp]
\centering
\includegraphics[width=\textwidth]{figure_6_forest_car_parts.pdf}
\caption{Forest plot for Car Parts (GIFT-Eval).  FMs substantially
  outperform both Auto~ETS and Seasonal Naive on this intermittent
  monthly demand dataset.}
\label{fig:forest_car_parts}
\end{figure}

\subsection{Sensitivity analyses}
\label{subsec:sensitivity}\label{sec:heterogeneity}

Table~\ref{tab:sensitivity} reports five sensitivity slices of the
$k = 420$ pool.  The results are robust: the sign is consistently
negative (FM better) across all slices except WAPE-only.  The
statistical significance varies with pool size and heterogeneity.

\begin{table}[htbp]
\centering
\caption{Sensitivity analyses on the model-level meta-regression.}
\label{tab:sensitivity}
\begin{tabular}{lrrlrl}
\toprule
Sensitivity & $k$ & $\hat{\Delta}$ & 95\,\% CI & $p$ & Note \\
\midrule
Full pool         & 420 & $-0.132$ & $[-0.234,\; -0.030]$ & 0.011 & Primary \\
Excl-M5           & 363 & $-0.126$ & $[-0.238,\; -0.015]$ & 0.026 & M5 not driving result \\
\WAPE{}-only      & 164 & $+0.097$ & $[-0.086,\; +0.279]$ & 0.296 & Null on point forecast \\
SQL-only          & 128 & $-0.345$ & $[-0.442,\; -0.248]$ & $< .0001$ & Strongest signal \\
Source B only     &  36 & $-0.229$ & $[-0.569,\; +0.110]$ & 0.179 & Low power ($k = 36$) \\
Source A only     & 384 & $-0.101$ & $[-0.206,\; +0.005]$ & 0.063 & Marginal \\
\bottomrule
\end{tabular}
\end{table}

\paragraph{Key observations.}
\begin{enumerate}
  \item \textbf{M5 is not driving the result.}  Excluding M5
    ($k = 363$) barely changes the estimate ($-0.126$ vs $-0.132$)
    and remains significant ($p = 0.026$).  This is a major change from
    the original $k = 10$ analysis, where the entire FM advantage came
    from M5's per-series \WAPE{} artefact.

  \item \textbf{SQL vs \WAPE{} is the dominant moderator.}  The SQL-only
    slice is the strongest finding ($p < .0001$); the WAPE-only slice is
    null.  This is not a dataset effect --- it persists within the same
    datasets and models.

  \item \textbf{Source B alone is underpowered.}  The 36~Source~B pairs
    show the right direction ($-0.229$) but are insufficient for
    significance at the 5\% level with only 3~dataset clusters.
    Source~A (fev-bench + GIFT-Eval) provides the statistical mass.
\end{enumerate}

\subsection{Heterogeneity diagnostics}
\label{subsec:heterogeneity-diagnostics}

\begin{table}[htbp]
\centering
\caption{Variance decomposition.
  $\hat{\sigma}^2_1$ is between-dataset; $\hat{\sigma}^2_2$ is
  between-paper within dataset.}
\label{tab:t62}
\begin{tabular}{lrrrrr}
\toprule
Model & $k$ & $Q$ & $\hat{\sigma}^2_1$ & $\hat{\sigma}^2_2$ & Note \\
\midrule
Full pool     & 420 & 115{,}285 & 0.012 & 0.031 & Primary \\
Excl-M5       & 363 &  70{,}795 & 0.019 & 0.013 & \\
SQL-only      & 128 &   6{,}043 & 0.008 & 0.029 & Lowest $Q$ \\
\WAPE{}-only  & 164 &  63{,}770 & 0.035 & 0.109 & Highest $\hat{\sigma}^2$ \\
\bottomrule
\end{tabular}
\end{table}

The high $Q$ in the full pool is driven by the metric split: mixing SQL
and \WAPE{} pairs forces the model to average across two effect
directions.  Within the SQL-only slice, $Q$ drops by an order of
magnitude and the variance components are moderate, indicating that
conditional on metric choice, the FM advantage is relatively
homogeneous across datasets and models.

\subsection{Cost-accuracy Pareto frontier}
\label{subsec:pareto}\label{sec:cost-accuracy}

The cost axis is the second research question.  Our baselines (\S5.2)
cost \$2.05 across three datasets on Azure CPU; Figure~\ref{fig:pareto_hw}
plots the per-dataset best accuracy against per-dataset cost on consumer
hardware (M-series MacBook, marginal electricity at \$0.006/hr).  The
expanded FM set (\S5.4) populates four FM points per dataset; the
fev-bench extraction adds cross-paper accuracy anchors.

\begin{figure}[htbp]
\centering
\includegraphics[width=\textwidth]{figure_6_4_pareto_consumer_hw.pdf}
\caption{Cost-accuracy Pareto frontier on consumer hardware.
  X-axis: $\log_{10}$ USD marginal electricity cost.
  Y-axis: best \WAPE{}/\WRMSSE{}.  The four Source~B FMs span the
  9\,M--120\,M parameter range.}
\label{fig:pareto_hw}
\end{figure}

\begin{figure}[htbp]
\centering
\includegraphics[width=\textwidth]{figure_6_5_pareto_cloud_gpu.pdf}
\caption{Cost-accuracy Pareto frontier re-costed on cloud GPU
  (V100 at \$1.24/hr).}
\label{fig:pareto_gpu}
\end{figure}

\subsection{Moderator outcomes}
\label{subsec:moderator-outcomes}\label{sec:moderator-results}\label{sec:rq2-moderators}

With $k = 420$ pairs across 9~datasets and 11~FMs, the meta-regression
has sufficient power for moderator analysis.  Moderators tested:

\begin{enumerate}
  \item \textbf{Metric} (\texttt{metric\_name}; $F(2, 417) = 20{,}231$,
    $p < .0001$).  The dominant moderator.
    SQL coefficient is large and negative (FM wins on distributional
    accuracy); \WAPE{} coefficient is positive (FM slightly worse).  See
    \S\ref{subsec:metric-moderator} for full results.

  \item \textbf{Dataset} (\texttt{dataset\_norm}; $F(8, 411) = 0.44$,
    $p = 0.90$).  Not significant as a moderator once metric is
    accounted for.  The raw $Q$-statistic shows extreme heterogeneity
    ($p < .0001$), but this is driven by the metric split, not by
    dataset-level differences.

  \item \textbf{Horizon bucket} (\texttt{horizon\_bucket};
    $F(2, 417) = 0.49$, $p = 0.61$).  No
    significant horizon effect on $\Delta$.  The FM advantage does not
    systematically increase or decrease with forecast horizon in this
    pool.

  \item \textbf{Model scale} (\texttt{model\_scale}, $\log_{10}$
    params; $F(1, 418) = 1{,}371$, $p < .0001$).  Coefficient is
    positive ($+0.087$ per $\log_{10}$ unit), meaning larger FMs show
    \emph{smaller} advantage.  However, this is confounded: model
    identity is correlated with which benchmark papers (and therefore
    which metrics and datasets) the model appears in.  Within Source~B,
    where protocol is held constant, the smallest FM (Moirai~2.0-Small,
    11\,M) outperforms the largest (Chronos-2, 120\,M) on M5, supporting
    \citet{shchur2025fevbench}'s finding that model scale matters less
    than architecture for retail forecasting.

  \item \textbf{Baseline family} (\texttt{baseline\_family};
    $F(1, 418) = 8.47$, $p = 0.004$).  ML\_TREE
    vs STATS.  The FM advantage is larger when the baseline is
    statistical (Auto~ETS, Seasonal Naive; $\hat{\Delta} = -0.367$)
    than when it is a gradient-boosted tree ($\hat{\Delta} = -0.079$),
    as expected: tree baselines are stronger competitors.
\end{enumerate}

All moderator specifications and outputs are in
\texttt{analysis/meta\_regression.R} and the corresponding
\texttt{analysis/figures/moderator\_*.txt} files.
